\documentclass[twocolumn]{article}

\usepackage[utf8]{inputenc}
\usepackage[margin=1.25in]{geometry}
\usepackage{authblk}
\usepackage{setspace}
\usepackage{graphicx}
\graphicspath{{./figures/}}
\usepackage{subcaption}
\usepackage{amsmath}
\usepackage{booktabs}
\usepackage{hyperref}

\setlength{\columnsep}{24pt}
\usepackage[switch]{lineno}
\setlength\linenumbersep{8pt}

\usepackage[
  style=nejm,
  citestyle=numeric-comp,
  sorting=none
]{biblatex}
\addbibresource{references.bib}

\title{Análisis de la huella de carbono en fine-tuning de modelos de lenguaje: implicaciones energéticas y ambientales}
\author[1]{Daniel Armas Ramírez}
\author[2]{Ian Josué López Sánchez}
\affil[1]{Escuela Superior de Cómputo, Instituto Politécnico Nacional, Ciudad de México, México.}
\affil[2]{Escuela Superior de Cómputo, Instituto Politécnico Nacional, Ciudad de México, México.}
\date{Junio de 2026}

\setstretch{1.15}

\begin{document}
\twocolumn[
\begin{@twocolumnfalse}

\maketitle
\begin{abstract}
The growing deployment of Large Language Models (LLMs) has intensified the energy demand of data centers, whose global consumption reached 460~TWh in 2022, representing two percent of worldwide electricity use. Fine-tuning, a transfer-learning technique that adapts pretrained models to specific domains, has emerged as an efficient alternative to training from scratch; however, its carbon footprint remains insufficiently quantified, particularly in regions such as Mexico. This work evaluates the carbon intensity of fine-tuning open-source LLMs ranging from 7 to 20 billion parameters, belonging to the LLaMA~3, Qwen~3, Gemma~3, BioMistral, and Ministral families, over the BioInstruct biomedical instruction dataset. The Software Carbon Intensity (SCI) methodology from the Green Software Foundation was applied, measuring GPU, CPU, and RAM consumption with CodeCarbon, Zeus, and Prometheus with Grafana, on an on-premise server with AMD Instinct MI210 GPUs and on Lightning~AI with NVIDIA A100 and L4 GPUs. The quality of the fine-tuned models was assessed using BERTScore with BioBERT as the reference model. Results indicate that model size does not linearly correlate with output quality; LLaMA~3 8B with LoRA achieved the best performance-to-environmental-cost ratio, while BioMistral 7B yielded the most sustainable configuration overall. The Mexican national electricity system emission factor (0.444~tCO$_2$e/MWh) was used to contextualize the environmental impact in the region.
\end{abstract}
\vspace{1em}
\end{@twocolumnfalse}
]

\section{Introducción}

Los modelos grandes de lenguaje (LLM), basados en la arquitectura transformer y sus mecanismos de autoatención \cite{googleTransformers}, han transformado el procesamiento del lenguaje natural. Sin embargo, su operación demanda cantidades significativas de energía: en 2022 los centros de datos consumieron 460~TWh a nivel global, equivalente al dos por ciento de la demanda eléctrica mundial \cite{ieaelectricity2024}, y se estima que una consulta a ChatGPT produce 4.32~g de CO$_2$ \cite{xia2024understandingperformanceestimatingcost}, mientras que el entrenamiento completo de un LLM puede emitir hasta 550 toneladas de CO$_2$ equivalente \cite{dodge2022measuringcarbonintensityai}.

Ante este panorama, el fine-tuning surge como alternativa al reutilizar pesos de modelos preentrenados y adaptarlos mediante aprendizaje supervisado \cite{ibmfinetuning}. Técnicas como PEFT y LoRA reducen drásticamente los parámetros entrenables al introducir matrices de bajo rango, disminuyendo memoria y tiempo de cómputo \cite{Lora,PEFTHugging}. No obstante, la cuantificación de su huella de carbono sigue siendo limitada, y prácticamente inexistente para el contexto energético mexicano.

Por lo anterior, este estudio evalúa la huella de carbono del fine-tuning de LLMs de código abierto de 7 a 20 mil millones de parámetros, empleando la metodología Software Carbon Intensity (SCI) \cite{greensoftwarefoundation_sci2025} y el factor de emisión del sistema eléctrico nacional (FESEN), comparando arquitecturas densas y de mezcla de expertos, con el fin de proponer estrategias hacia una IA más sostenible.

\section{Metodología}

Se adoptó la metodología SCI de la Green Software Foundation \cite{greensoftwarefoundation_sci2025}, que calcula las emisiones de un sistema de software como $SCI = ((E \times I) + M)$ por $R$, donde $E$ es la energía consumida en MWh, $I$ la intensidad de carbono de la red eléctrica, $M$ las emisiones incorporadas del hardware y $R$ la unidad funcional, definida en este trabajo como un experimento completo de fine-tuning. Los componentes monitoreados fueron GPU, CPU y RAM, mientras que la refrigeración e infraestructura de red se consideraron mediante el factor de eficacia del uso de energía (PUE) del servidor.

Los experimentos se ejecutaron en dos entornos: un servidor con AlmaLinux, dos GPUs AMD Instinct MI210, un CPU AMD EPYC y 384~GB de RAM orquestado con Podman; y estudios en Lightning~AI \cite{lightning2023} con GPUs NVIDIA A100 y L4. En ambos se realizó fine-tuning supervisado (SFT) \cite{trlsft2024} con LoRA \cite{Lora} y QLoRA \cite{qlora2023} sobre modelos de las familias LLaMA~3, Qwen~3, Gemma~3, BioMistral y Ministral \cite{yang2025qwen3technicalreport,gemma3report2025,mistral7b,ministral3}, utilizando el conjunto de datos biomédico BioInstruct \cite{Tran2024Bioinstruct}.

Para la medición energética se emplearon tres herramientas complementarias: CodeCarbon \cite{benoit_courty_2024_11171501} para la estimación agregada por experimento, Zeus \cite{measuringgpuzeus} para lecturas por fase en GPUs NVIDIA, y Prometheus con Grafana junto al AMD Device Metrics Exporter para la telemetría continua on-premise. Los contadores RAPL proporcionaron el consumo de CPU y RAM, mientras que AMD~SMI y NVML las lecturas de GPU. La calidad de los modelos ajustados se evaluó con BERTScore \cite{bertscore} usando \texttt{dmis-lab/biobert-v1.1} con parámetros fijos de 300 muestras y 256 tokens máximos de generación.

Finalmente, las emisiones operacionales se obtuvieron como $O = E \times I$, con $I$ igual al FESEN 2024 de 0.444~tCO$_2$e/MWh, y las emisiones incorporadas como $M = TE \times TS \times RS$, donde $TE$ son las emisiones del ciclo de vida del hardware, $TS$ la fracción temporal y $RS$ la fracción de recursos asignados. El puntaje SCI se calculó sumando ambos componentes por unidad funcional.

\section{Desarrollo}

El \textit{fine-tuning} se ejecutó mediante un pipeline modular de cinco fases secuenciales: carga y formateo del conjunto de datos, cuantización del modelo base, añadidura de adaptadores LoRA \cite{Lora}, entrenamiento supervisado y almacenamiento de adaptadores. Toda la ejecución se orquestó en contenedores Podman sobre el servidor mediante un archivo \texttt{docker-compose.yml} que levanta los servicios de Prometheus, Grafana, Node Exporter, el AMD Device Metrics Exporter y un contenedor de PyTorch con ROCm. El conjunto de datos BioInstruct \cite{Tran2024Bioinstruct} se reformatea al esquema de chat esperado por cada tokenizador mediante \texttt{apply\_chat\_template()}; el caso de gpt-oss \cite{openai2025gptoss120bgptoss20bmodel} requiere adicionalmente la conversión al formato Harmony \cite{openai2025harmony} estandarizado con Unsloth \cite{unslothgptoss2025}.

Las familias evaluadas fueron Mistral 7B \cite{mistral7b}, Ministral~3 14B \cite{ministral3}, gpt-oss 20B, Qwen~3 8B y 14B \cite{yang2025qwen3technicalreport}, LLaMA~3 8B y Gemma~3 12B \cite{gemma3report2025} quienes comparten una base Transformers, salvo gpt-oss que sigue un enfoque de mezcla de expertos. Los pesos se cargan en formato NormalFloat de 4 bits (NF4) con doble cuantización. Posteriormente, con el uso de PEFT \cite{PEFTHugging} se congela los pesos base, convierte las capas de normalización a FP32 y habilita el punto de control de gradientes; los adaptadores LoRA se entrenan con la clase \texttt{SFTTrainer} de TRL \cite{trlsft2024}, y al concluir sólo se mantienen los adaptadores y el tokenizador, lo que genera puntos de control.

Como contingencia frente a la saturación temporal del servidor por procesos concurrentes externos, parte de los experimentos se migraron a Lightning~AI \cite{lightning2023} con una GPU NVIDIA A100 de 80~GB, lo que permitió paralelizar la ejecución y concluir la fase experimental dentro de los tiempos planeados sin desviarse de la configuración descrita.

\section{Conclusiones}

Los resultados confirman que el tamaño del modelo no correlaciona con la calidad de generación ni con la eficiencia ambiental. LLaMA~3.1 8B con LoRA logró el mejor equilibrio rendimiento-costo ($F1 = 0.763$, $SCI = 140.59$~gCO$_2$eq/R), mientras que BioMistral 7B \cite{mistral7b} obtuvo el menor puntaje SCI del estudio (102.59~gCO$_2$eq/R) con un F1 de 0.684, posicionándose como la configuración más sostenible. En contraste, Gemma~3 12B \cite{gemma3report2025} con cuantización alcanzó el SCI más elevado (hasta 415.73~gCO$_2$eq/R) sin mejora proporcional en calidad, y Ministral~3 14B \cite{ministral3} presentó degradación por la combinación de LoRA y cuantización aplicada sobre un modelo multimodal evaluado en tareas exclusivamente textuales.

Estos hallazgos refuerzan que la selección de la estrategia de fine-tuning \cite{ibmfinetuning} y la arquitectura base tienen un impacto más determinante sobre el SCI que el número de parámetros. Los modelos densos de menor escala (7B--8B) ajustados con LoRA y cargados sin fusionar el adaptador superan a alternativas de mayor tamaño tanto en calidad textual como en eficiencia ambiental, alineándose con la crítica reciente al paradigma \textit{bigger-is-better} en IA \cite{10.1145/3715275.3732006} y subrayando la viabilidad de configuraciones sostenibles bajo las condiciones energéticas mexicanas.

\printbibliography

\end{document}